% Source: https://artofproblemsolving.com/community/c3h1881916_how_are_you_supposed_to_actually_solve_this
% Stable identity: aops-topic-1881916-q01
\textbf{No rational roots.} By the rational root theorem a rational root would be one of $\pm1,\pm2,\pm4$, and none of these works ($f(-1)=2$, $f(-2)=-2$, $f(-4)=-40$, and $f(x)>0$ for $x>0$). So we need a formula.

\textbf{Remove the $x^2$ term.} Put $x=\dfrac{z-2}{3}$, i.e.\ $z=3x+2$. Then $27\left(x^3+2x^2+3x+4\right)=z^3+15z+70$ (expand $(z-2)^3+6(z-2)^2+27(z-2)+108$). So we must solve
\[
z^3+15z+70=0 .
\]
\textbf{Exactly one real root.} The function $z\mapsto z^3+15z+70$ is strictly increasing (a sum of increasing functions), so there is exactly one real $z$, hence exactly one real $x$.

\textbf{Cardano's method.} Look for $z=u+v$ with $uv=-5$. Then $z^3=u^3+v^3+3uv(u+v)=u^3+v^3-15z$, so the equation becomes $u^3+v^3=-70$. Together with $u^3v^3=-125$, the numbers $u^3$ and $v^3$ are the roots of $s^2+70s-125=0$, i.e.\ $s=-35\pm\sqrt{1350}=-35\pm15\sqrt6$. Taking real cube roots (their product is $\sqrt[3]{(15\sqrt6-35)(-15\sqrt6-35)}=\sqrt[3]{-125}=-5$, as required),
\[
z=\sqrt[3]{15\sqrt6-35}-\sqrt[3]{15\sqrt6+35}.
\]
Therefore the only real solution is
\[
\boxed{x=\dfrac{-2+\sqrt[3]{15\sqrt6-35}-\sqrt[3]{15\sqrt6+35}}{3}}\approx-1.6506 .
\]
(Consistently, $f(-2)<0<f(-1)$; the other two roots are non-real.)
